% Part: normal-modal-logic
% Chapter: filtrations
% Section: S5-fmp

\documentclass[../../../include/open-logic-section]{subfiles}

\begin{document}

\olfileid{nml}{fil}{fmp}

\olsection{\Log{K} lan \Log{S5} Nduweni Sipat Model Winates}

\begin{defn}
  Sistem $\Sigma$ saka logika modal diarani nduweni \emph{sipat model
    winates} yen saben !!a{formula}~$!A$ kang bener ing sawijining
  jagad ing model saka $\Sigma$, $!A$ uga bener ing sawijining jagad ing
  model \emph{winates} saka~$\Sigma$.
\end{defn}

\begin{prop}\ollabel{prop:K-fmp}
  \Log{K} nduweni sipat model winates.
\end{prop}

\begin{proof}
  \Log{K} iku himpunan !!{formula}s kang valid, yaiku sembarang model
  iku model saka~\Log{K}. Miturut \olref[fil]{thm:filtrations}, yen
  $\mSat{M}{!A}[w]$, mula $\mSat{M^*}{!A}[{[w]}]$ kanggo sembarang
  filtrasi saka~$\mModel{M}$ liwat himpunan $\Gamma$ sub-!!{formula}s
  saka~$!A$. Saben !!{formula} mung nduweni sub-!!{formula}s kang
  cacahé winates, mula $\Gamma$ winates. Miturut
  \olref[fin]{prop:filt-are-finite}, $\card{W^*} \le 2^n$, ing ngendi
  $n$ iku cacah !!{formula}s ing~$\Gamma$. Lan amarga \Log{K} ora
  mbatesi model, $\mModel{M^*}$ iku model \Log{K}.
\end{proof}

Kanggo nuduhake yen logika~\Log{L} nduweni sipat model winates liwat
filtrasi, filtrasi saka model \Log{L} kudu uga dadi model \Log{L}.
Asring iki mbutuhake bukti kang ora prasaja, lan ora saben filtrasi
ngasilake model \Log{L}. Nanging kanggo model universal, sipat iki
tetep lumaku.

\begin{prop}\ollabel{prop:univ-fin}
  Upamane $\mClass{U}$ kelas model universal (delengen
  \olref[frd][es5]{prop:S5=univ}) lan $\mClass{U}_\mathrm{Fin}$ kelas
  kabeh model universal winates. Banjur saben !!{formula}~$!A$ valid
  ing $\mClass{U}$ yen lan mung yen valid ing
  $\mClass{U}_\mathrm{Fin}$.
\end{prop}

\begin{proof}
  Model universal winates kalebu model universal, mula arah kiwa
  menyang tengen langsung cetha. Kanggo arah tengen menyang kiwa,
  upamane~$!A$ salah ing sawijining jagad $w$ ing model universal
  $\mModel{M}$. Jupuken $\Gamma$ kang ngemot $!A$ lan kabeh
  subformulane; cetha yen $\Gamma$ winates. Jupuken filtrasi
  $\mModel{M^*}$ saka $\mModel{M}$; banjur $\mModel{M^*}$ winates
  miturut \olref[fin]{prop:filt-are-finite}, lan miturut
  \olref[fil]{thm:filtrations}, $!A$ salah ing $[w]$ ing
  $\mModel{M^*}$. Isih kudu digatekake yen $\mModel{M^*}$ uga
  universal: kanggo $u$ lan $v$, miturut hipotesis $Ruv$ lan miturut
  \olref[fil]{defn:filtration}\olref[fil]{defn:filtration-R},
  uga $R^*[u][v]$.
\end{proof}

\begin{cor}\ollabel{cor:S5fmp}
  \Log{S5} nduweni sipat model winates.
\end{cor}

\begin{proof}
  Miturut \olref[frd][es5]{prop:S5=univ}, yen $!A$ bener ing sawijining
  jagad ing model refleksif lan Euklidean, mula bener ing sawijining
  jagad ing model universal. Miturut \olref{prop:univ-fin}, formula iku
  bener ing sawijining jagad ing model universal winates (yaiku
  filtrasi model liwat himpunan sub-!!{formula}s saka~$!A$). Saben
  model universal uga refleksif lan Euklidean; mula $!A$ bener ing
  sawijining jagad ing model refleksif lan Euklidean kang winates.
\end{proof}

\begin{prob}
  Tuduhna yen saben filtrasi saka model serial utawa refleksif uga
  serial utawa refleksif (miturut urutane).
\end{prob}

\begin{prob}
  Temokna filtrasi kang ora simetris (ora transitif, ora Euklidean)
  saka model simetris (transitif, Euklidean).
\end{prob}

\end{document}
